\clearpage\begingroup
\ArithmeticHead{13}{R}{III. ÜBUNG DES DIVISIONSVERFAHRENS AN KETTENBRÜCHEN.}
\renewcommand{\thefootnote}{\arabic{footnote})}\setcounter{footnote}{0}
\newcommand{\MDown}[2]{\mathord{\raisebox{#1}{$\displaystyle #2$}}}
\begin{center}III.\\[3pt]{\large\textls[100]{Übung des Divisionsverfahrens an Kettenbrüchen.}}\end{center}
\noindent\hspace*{\parindent}Eine andere Anfangsübung ist die Bestimmung von $\sqrt{5}$\footnote{Nachlass, Kapsel 44.}:
\[
\sqrt{5}=3-\frac{1}{4}-\frac{6}{13}+\frac{6}{29}+\frac{34}{211}-\frac{177}{421}=\frac{299537289}{133957148}.
\]
\noindent Hier benutzt er bereits die Kettenbruchentwicklung für $\sqrt{5}$, und zwar den
12. Näherungsbruch, zerlegt den Nenner in die Faktoren $4\mathord{.}13\mathord{.}29\mathord{.}211\mathord{.}421$,
macht dann in der oben geschilderten Weise die Partialbruchzerlegung und
erhält so das rechts von $\sqrt{5}$ stehende Aggregat. Auch den 8. Näherungswert
derselben Kettenbruchentwicklung gibt er darunter an:
\[
\sqrt{5}=\frac{930249}{416020}\qquad 416020=4\mathord{.}5\mathord{.}11\mathord{.}31\mathord{.}61.
\]
\clearpage\SourceLeaf{569}\ArithmeticHead{14}{L}{PHILIPP MAENNCHEN, GAUSS ALS ZAHLENRECHNER.}
\setcounter{footnote}{0}
\noindent Eine vollständig durchgeführte Übungsaufgabe dieser Art finden wir \textsc{Leiste},
S. 50. Es handelt sich um $\sqrt{17}=4+\frac{1}{8}\MDown{-.4ex}{+\frac{1}{8}}\MDown{-.8ex}{+\ldots\ \text{in inf.}}$
Man beachte wieder die praktische Anordnung.
\begingroup\small\setlength{\arraycolsep}{3pt}
\[
\begin{array}{ccccccccc}
\scriptstyle1&\scriptstyle2&\scriptstyle3&\scriptstyle4&\scriptstyle5&\scriptstyle6&\scriptstyle7&\scriptstyle8&\scriptstyle9\\
\frac{4}{1}&\frac{33}{8}&\frac{268}{65}&\frac{2177}{528}&\frac{17684}{4289}&\frac{143649}{34840}&\frac{1166876}{283009}&\frac{9478657}{2298912}&\frac{}{65\mathord{.}287297}\\
&&&&\text{pr.}&&\text{pr.}&&
\end{array}
\]
\endgroup
\[
\frac{9478657}{2298912}=\frac{5}{7}+\frac{2}{3}+\frac{4}{11}+\frac{1}{32}+\frac{730}{311}
\]
\begingroup\small
\[
\begin{array}{r@{\;}l}
\frac{730}{311}={}&2{,}347266881028938906752411575562\\[5pt]
\frac{1}{32}={}&\phantom{0{,}0}3125\\[5pt]
\text{Cetera}={}&1{,}744588744588744588744588744588\\\hline
&4{,}12310562561768349549
\end{array}\text{\textsuperscript{1)}}.
\]
\endgroup
\setcounter{footnote}{1}\footnotetext{Er addiert zunächst $\frac{2}{3}+\frac{4}{11}=1+\frac{1}{33}=1{,}030303$, und dieses Ergebnis zählt er zu $\frac{5}{7}=0{,}714285714285$ in der Dezimalbruchtabelle.}
\par\noindent\hspace*{\parindent}Man sieht, dass der 5. Näherungswert bereits auf seine Brauchbarkeit
untersucht worden ist; allein der Nenner erwies sich als Primzahl (pr.).
Ebenso war es mit dem Nenner des 7. Näherungswertes. Der 6. war zwar
brauchbar; denn der Nenner ist $40\mathord{.}13\mathord{.}67$; aber er wurde wohl einstweilen
zurückgestellt. Ein Näherungswert für eine Quadratwurzel scheint unbrauchbar zu sein, wenn der Nenner nicht in Primfaktoren unter 1000 zerlegbar ist,
da in diesem Falle die Dezimalbruchtabelle nicht benutzt werden kann. Man
kann aber, wenn der Zähler in dieser Weise zerlegbar ist,
\[
\frac{1}{N}\approx\frac{1}{\sqrt{17}}=\frac{1}{17}\sqrt{17}
\]
\noindent berechnen und dann mit 17, oder allgemein mit dem Radikanden multiplizieren. Solche Versuche, den Zähler zu zerlegen, finden sich gelegentlich
bei \textsc{Gauss}.

Über diese Stufe erhebt sich nun \textsc{Gauss} durch die Versuche, die \textls[100]{Genauigkeit zu verschärfen}. In Kapsel 44 des Nachlasses findet sich ein
Blatt mit dem folgenden Beispiel:
\[
\begin{aligned}
\sqrt{2}&=\frac{275807^2+2\mathord{.}195025^2}{2\mathord{.}195025\mathord{.}275807}=\frac{152139002499}{107578520350}\quad\text{quam proxime}\\
&=\frac{7}{50}+\frac{1}{7}+\frac{2}{29}+\frac{37}{41}+\frac{134}{269}+\frac{636}{961}-1.
\end{aligned}
\]
\clearpage\SourceLeaf{570}\ArithmeticHead{15}{R}{III. ÜBUNG DES DIVISIONSVERFAHRENS AN KETTENBRÜCHEN.}
\noindent Der Gedankengang ist vermutlich der folgende gewesen: Für
\[
\sqrt{2}=1+\frac{1}{2}\MDown{-.4ex}{+\frac{1}{2}}\MDown{-.8ex}{+\ldots\ \text{in inf.}}
\]
\noindent ist der 14. Näherungswert $\frac{275807}{195025}$. Nun sei $\frac{u}{v}$ ein Näherungswert für $\sqrt{n}$ und
zwar $\frac{u}{v}>\sqrt{n}$; dann ist $\frac{u}{v\sqrt{n}}>1$ und $\frac{v\sqrt{n}}{u}<1$. Alsdann ist mit noch viel
geringerer Abweichung die Gleichung erfüllt:
\[
\frac{u}{v\sqrt{n}}+\frac{v\sqrt{n}}{u}=2,
\]
\noindent woraus folgt:
\[
\frac{u}{v}+\frac{nv}{u}=2\sqrt{n},
\]
\noindent und hieraus:
\[
\frac{u^2+nv^2}{2uv}=\sqrt{n},
\]
\par\noindent\hspace*{\parindent}Noch wahrscheinlicher ist allerdings der folgende Gedankengang, der ein
besseres Urteil über den Grad der erreichten Genauigkeit ermöglicht:
\[
\begin{aligned}
&\sqrt{n}\approx\frac{u}{v};\quad v^2\mathord{.}n=u^2\pm\varepsilon\\
&\hspace{3em}n=\frac{u^2\pm\varepsilon}{v^2}=\frac{u^2}{v^2}\left(1\pm\frac{\varepsilon}{u^2}\right)\\
&\hspace{2em}\sqrt{n}=\frac{u}{v}\left(1\pm\frac{\varepsilon}{u^2}\right)^{\frac12}=\frac{u}{v}\left(1\pm\frac{\varepsilon}{2u^2}+\ldots\right)=\frac{u}{v}\pm\frac{\varepsilon}{2uv}\\
&\hspace{2em}\sqrt{n}=\frac{2u^2\pm\varepsilon}{2uv},
\end{aligned}
\]
\noindent und da $u^2\pm\varepsilon=v^2\mathord{.}n$ ist,
\[
\sqrt{n}=\frac{u^2+nv^2}{2uv}.
\]
\par\noindent\hspace*{\parindent}Auf der nächsten Stufe zeigt \textsc{Gauss} bereits eine Rechengewandtheit, der
nur schwer zu folgen ist. Die Kettenbruchentwicklung liefert ihm für $\sqrt{15}$:
\begingroup\small
\[
\begin{aligned}
\text{1)}\quad\sqrt{15}\ \text{quam proxime}&=\frac{457470751}{118118440}=4+\frac{3}{8}+\frac{2}{5}-\frac{2}{11}-\frac{9}{19}-\frac{20}{71}+\frac{7}{199}\\[4pt]
\text{2)}\quad\sqrt{15}\ \text{quam proxime}&=\frac{28355806081}{7321437648}=5+\frac{5}{16}-\frac{1}{7}+\frac{1}{9}-\frac{9}{23}-\frac{11}{31}-\frac{21}{61}-\frac{53}{167}.
\end{aligned}
\]
\endgroup
\noindent Nun schreibt er Ergebnis und Interpolationen:
\begingroup\small\setlength{\arraycolsep}{2pt}
\[
\begin{array}{r@{\;}r@{\;}l}
\text{zu 1)}&3{,}872983\,346207\,416894\,432401\,917939\,315825\,708500\,721817\,8635&\text{etc.}\\
&-\quad9\,253136\,518156\,916248\,036269\,795441\,7885&\text{etc.}\\
&416885\,179265\,399782\,399577\,672230\,926376\,0750&\text{etc.}\\
&+\qquad33\,160690\,778915\,5160&\text{etc.}\\[-2pt]
&\overline{399610\,832921\,705291\,5910}&
\end{array}
\]
\endgroup
\clearpage\SourceLeaf{571}\ArithmeticHead{16}{L}{PHILIPP MAENNCHEN, GAUSS ALS ZAHLENRECHNER.}
\begingroup\small\setlength{\arraycolsep}{2pt}
\[
\begin{array}{r@{\;}r@{\;}l}
\text{zu 2)}&3{,}872983\,346207\,416885\,181673\,816530\,193032\,929862\,629624\,3505&\text{etc.}\\
&-\quad2408\,416747\,793422\,096943\,170845\,6075&\text{etc.}\\
&179265\,399782\,399610\,832919\,458778\,7430&\text{etc.}\\
&+\qquad2\,246512\,8480&\text{etc.}\\[-2pt]
&\overline{832921\,705291\,5910}&\text{etc.}
\end{array}
\]
\endgroup
\par\noindent\hspace*{\parindent}Man sieht: Der erste Näherungswert liefert $\sqrt{15}$ auf 16 Dezimalstellen;
der zweite auf 19 Stellen; die zweimalige Interpolation auf mehr als 50
Stellen! Eine Andeutung der Methode war nirgends zu finden, wir wollen
daher versuchen, den Weg aufzudecken, den \textsc{Gauss} gegangen sein könnte.

Bezeichnet man den von \textsc{Gauss} unter 1) angegebenen, aus der Kettenbruchentwicklung
\[
\sqrt{15}=4-\frac{1}{8}\MDown{-.4ex}{-\frac{1}{8}}\MDown{-.8ex}{-\ \text{in inf.}}
\]
\noindent gewonnenen Wert mit $\frac{Z}{N}$, so besteht auf Grund dieser Entwicklung die
Gleichung
\[
15N^2=Z^2-1,
\]
\noindent also
\[
\begin{aligned}
15&=\frac{Z^2-1}{N^2}=\frac{Z^2}{N^2}\left(1-\frac{1}{Z^2}\right);\\
\sqrt{15}&=\frac{Z}{N}\left(1-\frac{1}{Z^2}\right)^{\frac12}.
\end{aligned}
\]
\noindent Berücksichtigt man nur die drei ersten Glieder der Binomialreihe, so entsteht die Gleichung:
\[
\sqrt{15}=\frac{Z}{N}\left(1-\frac{1}{2}\mathord{.}\frac{1}{Z^2}-\frac{1}{8}\mathord{.}\frac{1}{Z^4}\right).
\]
\par\noindent\hspace*{\parindent}Diese Entwicklung war \textsc{Gauss} naturgemäss nicht »zierlich« genug. Denn
die Division durch $Z^2$, also hier durch $457470751^2$, war unvermeidlich, und
wir sehen ja, dass er bereits der Division durch $N$, also durch 118118440
aus dem Wege gegangen ist, indem er $\frac{Z}{N}$ in Teilbrüche zerlegt hat, die er
mit Hilfe seiner Tafel der reziproken Werte der Primzahlen leicht auf beliebig viele Stellen hinschreiben konnte. Zu diesem Zweck hatte er 118118440
in seine Primfaktoren zerlegt; und nun müsste er in gleicher Weise 457470751
zerlegen. Das wäre erstens umständlich, und zweitens wäre der Erfolg zweifelhaft.
\clearpage\SourceLeaf{572}\ArithmeticHead{17}{R}{III. ÜBUNG DES DIVISIONSVERFAHRENS AN KETTENBRÜCHEN.}
\noindent\hspace*{\parindent}Dagegen ist die Division durch $Z^2-1$ aussichtsreich; denn diese Zahl
ist nach unserer ersten Formel gleich $15N^2$, und $N$ ist ja bereits in Faktoren
zerlegt. Statt der Division durch $457470751^2$ wird daher die durch
\[
457470751^2-1=15\mathord{.}(8\mathord{.}5\mathord{.}11\mathord{.}19\mathord{.}71\mathord{.}199)^2
\]
\noindent gewählt.
\par\noindent\hspace*{\parindent}Nun ist
\[
\begin{gathered}
\frac{1}{Z^2-1}=\frac{\frac{1}{Z^2}}{1-\frac{1}{Z^2}}=\frac{1}{Z^2}\left(1+\frac{1}{Z^2}+\frac{1}{Z^4}+\ldots\right);\quad\text{also}\quad\frac{1}{Z^2}=\frac{1}{Z^2-1}-\frac{1}{Z^4}-\ldots\\[6pt]
\begin{aligned}
\sqrt{15}&=\frac{Z}{N}-\frac{1}{2}\mathord{.}\frac{Z}{N}\left(\frac{1}{Z^2}+\frac{1}{4}\mathord{.}\frac{1}{Z^4}\right)=\frac{Z}{N}-\frac{1}{2}\mathord{.}\frac{Z}{N}\mathord{.}\frac{1}{Z^2}\left(1+\frac{1}{4}\mathord{.}\frac{1}{Z^2}\right)\\
&=\frac{Z}{N}-\frac{1}{2}\mathord{.}\frac{Z}{N}\mathord{.}\frac{1}{Z^2-1}\left(1-\frac{3}{4}\mathord{.}\frac{1}{Z^2}\right).
\end{aligned}
\end{gathered}
\]
\noindent In dem Klammerausdruck setzt man für $\frac{1}{Z^2}$ den Wert $\frac{1}{Z^2-1}$, ohne die oben
abgegrenzte Genauigkeit zu beeinträchtigen. Es ergibt sich also:
\[
\sqrt{15}=\frac{Z}{N}-\left(\frac{Z}{N}\right)\mathord{.}\frac{1}{2}\mathord{.}\frac{1}{Z^2-1}+\left(\frac{Z}{N}\mathord{.}\frac{1}{2}\mathord{.}\frac{1}{Z^2-1}\right)\mathord{.}\frac{3}{4}\mathord{.}\frac{1}{Z^2-1}.
\]
\par\noindent\hspace*{\parindent}Mit dem Einklammern soll angedeutet werden, dass etwas bereits Berechnetes zur weiteren Rechnung benutzt wird, also zuerst der Näherungswert $\frac{Z}{N}$, dann die erste Interpolationsgrösse.
\par\noindent\hspace*{\parindent}Nach dieser theoretischen Auseinandersetzung soll nun die ganze Rechnung durchgeführt werden.
\begingroup\small\setlength{\arraycolsep}{3pt}\renewcommand{\arraystretch}{1.2}
\[
\begin{array}{r@{\;}l}
\frac{Z}{N}={}&3{,}872983\,346207\,416894\,432401\,917939\,315825\,708500\,721817\,8635\\
{}[:\;40]&0{,}096824\,583655\,185422\,360810\,047948\,482895\,642712\,518045\,4465\\
{}[:\;40]&0{,}002420\,614591\,379635\,559020\,251198\,712072\,391067\,812951\,1361\\
{}[:\;11]&0{,}000220\,055871\,943603\,232638\,204654\,428370\,217369\,801177\,3760\\
{}[:\;11]&0{,}000020\,005079\,267600\,293876\,200423\,129851\,837942\,709197\,9432\\
{}[:\;19]&0{,}000001\,052898\,908821\,068098\,747390\,691044\,833575\,932063\,0496\\
{}[:\;19]&0{,}000000\,055415\,732043\,214110\,460388\,983739\,201767\,154319\,1078\\
{}[:\;71]&0{,}000000\,000780\,503268\,214283\,245920\,971601\,960588\,269779\,1423\\
{}[:\;71]&0{,}000000\,000010\,993003\,777665\,961210\,154529\,605078\,708025\,0583\\
{}[:\;199]&0{,}000000\,000000\,055241\,225013\,396790\,000776\,530678\,787477\,5128\\
{}[:\;199]&0{,}000000\,000000\,000277\,594095\,544707\,487441\,088093\,863253\,6558\\
{}[:\;30]&0{,}000000\,000000\,000009\,253136\,518156\,916248\,036269\,795441\,7885.
\end{array}
\]
\endgroup
\par\smallskip\noindent{\small X\textsubscript{2} Abh. 6.}\hfill{\small 3}\par
\clearpage\SourceLeaf{573}\ArithmeticHead{18}{L}{PHILIPP MAENNCHEN, GAUSS ALS ZAHLENRECHNER.}
\noindent\hspace*{\parindent}Damit ist der erste Interpolationswert
\[
J_1=\frac{Z}{N}\mathord{.}\frac{1}{2}\mathord{.}\frac{1}{Z^2-1}=\frac{Z}{N}\mathord{.}\frac{1}{2}\mathord{.}\frac{1}{15}\mathord{.}\frac{1}{(8\mathord{.}5\mathord{.}11\mathord{.}19\mathord{.}71\mathord{.}199)^2}
\]
\noindent berechnet. Hieraus gewinnt man den zweiten, indem man den ersten mit\linebreak
$\frac{3}{4}\mathord{.}\frac{1}{15}\mathord{.}\frac{1}{(8\mathord{.}5\mathord{.}11\mathord{.}19\mathord{.}71\mathord{.}199)^2}$ multipliziert. Wir wollen auch diese Rechnung ausführen und dabei, wie es auch \textsc{Gauss} vielfach zu tun pflegte, auf die vorausgehenden Nullen keine Rücksicht nehmen.
\begingroup\small\setlength{\arraycolsep}{3pt}\renewcommand{\arraystretch}{1.2}
\[
\begin{array}{r@{\;}l}
J_1={}&.\quad.\quad.\quad9253136\,518156\,916248\,036269\,795441\,7885\\
{}[:\;40]&.\quad.\quad.\quad231328\,412953\,922906\,200906\,744886\,0447\\
{}[:\;40]&.\quad.\quad.\quad.\quad5783\,210323\,848072\,655022\,668622\,1511\\
{}[:\;11]&.\quad.\quad.\quad.\quad.\quad525\,746393\,077097\,514092\,969874\,7410\\
{}[:\;11]&.\quad.\quad.\quad.\quad.\quad47\,795126\,643372\,501281\,179079\,5219\\
{}[:\;19]&.\quad.\quad.\quad.\quad.\quad2\,515532\,981230\,131646\,377847\,3432\\
{}[:\;19]&.\quad.\quad.\quad.\quad.\quad.\quad132396\,472696\,322718\,440939\,3338\\
{}[:\;71]&.\quad.\quad.\quad.\quad.\quad.\quad.\quad1864\,739052\,060883\,358323\,0892\\
{}[:\;71]&.\quad.\quad.\quad.\quad.\quad.\quad.\quad.\quad26\,263930\,310716\,667018\,6350\\
{}[:\;199]&.\quad.\quad.\quad.\quad.\quad.\quad.\quad.\quad.\quad131979\,549300\,083753\,8624\\
{}[:\;199]&.\quad.\quad.\quad.\quad.\quad.\quad.\quad.\quad.\quad.\quad663\,213815\,578310\,3209\\
{}[:\;20]&.\quad.\quad.\quad.\quad.\quad.\quad.\quad.\quad.\quad.\quad33\,160690\,778915\,5160.
\end{array}
\]
\endgroup
\par\noindent\hspace*{\parindent}Damit haben wir den zweiten Interpolationswert $J_2$, der mit dem von
\textsc{Gauss} berechneten genau übereinstimmt. Wer ein besonders klares Bild von
der Gewandtheit, Sicherheit und Ausdauer des 16- bis 17jährigen \textsc{Gauss} erhalten will, der möge das obige Beispiel einmal selbst vollständig durchrechnen.
Bezüglich der rechnerischen Durchführung seien noch einige Bemerkungen
angebracht.

Die Darstellung ist auch hier wieder, wie bei andern in dieser Arbeit
behandelten Beispielen, so, dass die Stellenzahl beliebig vergrössert werden
kann, ohne dass früher ausgeführte Rechnungen wiederholt werden müssen.
Denn wie man zu einer beliebigen Stelle eines Divisionsergebnisses unter Benutzung der vorhergehenden und der nachfolgenden den zu der betreffenden
Stelle gehörigen Rest ermittelt, das braucht wohl nicht besonders ausgeführt
zu werden. Der junge \textsc{Gauss} beherrschte dieses Verfahren sicherlich, und es
\clearpage\SourceLeaf{574}\ArithmeticHead{19}{R}{III. ÜBUNG DES DIVISIONSVERFAHRENS AN KETTENBRÜCHEN.}
\setcounter{footnote}{0}
\noindent kam ihm zweifellos u. a. bei der Anwendung seiner Dezimalbruchtabelle\footnote{Werke II, S. 411ff.} sehr
zustatten. Die Reste der Divisionen durch 40, 11, 19, 71, 199 hat \textsc{Gauss}
höchst wahrscheinlich »im Sinne behalten«, wenn es auch nicht ausgeschlossen
ist, dass er hie und da, um eine Ruhepause zu machen, den einen oder den
andern Rest auf einem Täfelchen notierte. Nun sind allerdings die bei unserm
Beispiel auftretenden Primzahlen hierzu ganz gut geeignet. Für 11 und 19
braucht dies keinen besonderen Nachweis; die Vielfachen von 71 sind leicht
zu bilden, und $199=200-1$ gibt die Möglichkeit, die Divisionsreste leicht
im Kopf zu berechnen. Die Durchführung ist also im \textsc{Gauss}schen Sinne
»zierlich«, wenn auch die Rechenfreudigkeit eines \textsc{Gauss} dazu gehört, bis auf
52 Dezimalstellen zu gehen.

Bei dem von \textsc{Gauss} unter 2) angegebenen Näherungswert
\[
\frac{Z'}{N'}=\frac{28355806081}{7321437648}=5+\frac{5}{16}-\frac{1}{7}+\frac{1}{9}-\frac{9}{23}-\frac{11}{31}-\frac{21}{61}-\frac{53}{167}
\]
\noindent gelten für $Z'$, $N'$, $J'_1$ und $J'_2$ die gleichen Beziehungen wie für $Z$, $N$, $J_1$, $J_2$.
Die Durchführung gestaltet sich scheinbar weniger angenehm, da die Division
durch 23 und noch mehr die durch 167 abschrecken.

Wir dürfen jedoch sicher sein, dass \textsc{Gauss} auch dieser Unannehmlichkeit
ausgewichen ist. So kann man z. B. die Ergebnisse
\[
3\mathord{.}23=70-1
\]
\noindent und
\[
3\mathord{.}167=500+1
\]
\noindent benutzen, um weit bequemer zu dividieren. Ob \textsc{Gauss} gerade diese Beziehungen benutzte, oder ob er andere aufgespürt hat, das lässt sich freilich
nicht feststellen; jedenfalls gehen bei Benutzung der beiden erwähnten Beziehungen die Rechenschwierigkeiten wieder auf ein bescheidenes Mass zurück.

Warum hat \textsc{Gauss} die Aufgabe zweimal gelöst? War es die Freude am
Rechnen, war es die Absicht, die Rechenfertigkeit durch Übung zu steigern,
diente die eine Rechnung als Kontrolle der andern, wollte er den Grad der
Genauigkeit rechnerisch ergründen? Vermutlich haben damals noch alle diese
Gründe mitgespielt. Und dann ist möglicherweise noch eine andere Ursache
dazugekommen. Von den beiden Werten $\frac{Z}{N}$ und $\frac{Z'}{N'}$ ist der zweite als späterer
Näherungswert in der Kettenbruchentwicklung der genauere, und die Frage
\par\smallskip\noindent\hfill{\small 3\textsuperscript{*}}\par
\clearpage\SourceLeaf{575}\ArithmeticHead{20}{L}{PHILIPP MAENNCHEN, GAUSS ALS ZAHLENRECHNER.}
\noindent liegt nahe, wie man mit Hilfe dieser beiden Näherungswerte dem wahren
Werte möglichst nahe kommen kann. Diese Frage in allgemeinerer Fassung
hat ihn bekanntlich intensiv beschäftigt und in der berühmten Methode der
kleinsten Quadrate ihren Abschluss gefunden.
\par\clearpage\endgroup
